\begin{table}

\caption{Lead-time sensitivity of the augmented-observations precision gain, hurricane-force realizations. \label{tab:appB_leadtime}}
\centering
\begin{adjustbox}{max width = \textwidth}
\begin{threeparttable}
\begin{tabular}[t]{crrrrr}
\toprule
Lead (days) & Obs. & $\bar{\sigma}$ (m/s) & $\delta$ (m/s) & SE & p\\
\midrule
1 & 144 & 7.18 & 0.063 & (0.071) & 0.373\\
2 & 140 & 7.96 & 0.144 & (0.048) & 0.003\\
3 & 124 & 9.11 & 0.258 & (0.095) & 0.007\\
\bottomrule
\end{tabular}
\begin{tablenotes}
\small
\item Each row reports the controls-table specification of Table \ref{tab:noaa-controls} estimated on the indicated lead-time subset of hurricane-force observations (wind $\geq$ 33 m/s) in the 2017--2020 continental-US AO denial archive (excluding Hurricane Maria). The coefficient $\delta$ is the average reduction in the ex-ante wind-speed forecast standard deviation attributable to the assimilation of dropsonde profiles (the ALL configuration of \citet{DitchekEA2023}); the precision gain is reported as $|\hat{\delta}|$, so larger positive values indicate a larger noise reduction from AO. County-clustered standard errors in parentheses.
\end{tablenotes}
\end{threeparttable}
\end{adjustbox}
\end{table}
